% ! TeX root = ../main.tex

\section{Evaluation}
\label{sec:orca:eval}

We evaluate ORCA by progressively enabling its capabilities.
As a passive telemetry collector, we compare runtime overhead (\S\ref{sec:orca:eval-trace-runtime}),
storage efficiency (\S\ref{sec:orca:eval-trace-size}), and analytics performance
(\S\ref{sec:orca:eval-trace-query}) against state-of-the-art tracers.
We then demonstrate online analytics (\S\ref{sec:orca:eval-insitu}) and control workflows
(\S\ref{sec:orca:eval-agent}) --- capabilities with no existing counterparts.

\parahead{Hardware Setup}
Experiments run on a 600-node research cluster with \qlogic{} Truescale fabric (Intel Xeon E5-2670,
16 cores/node, 64GB RAM, 40\,Gbps \qlogic{} 7340 NICs) and a 20-node Lustre filesystem with 3GB/s
aggregate bandwidth.

\parahead{Software Setup}
We use the Sedov Blast Wave 3D~\cite{Sedov1959:SimilarityDimensionalMethods} AMR code from the
Parthenon~\cite{Grete2023:Parthenon} and Phoebus~\cite{Barke2024:Phoebus} frameworks at 512--4096
CPU ranks, with 16 ranks/node.
ORCA uses one aggregator node per 1024 ranks plus one controller node.
All experiments use MVAPICH2 v2.3.7~\cite{Panda2021:MVAPICH} with a tuned PSM~\cite{2025:PSM}
library for our fabric.
While simulation nodes use PSM, ORCA communication uses Mercury~\cite{Souma2013:MercuryRPC} over
verbs because of fabric restrictions.

This setup represents a strenuous observability workload: CPU codes make every cycle visible as
runtime overhead, the legacy fabric lacks QoS support and forces ORCA onto a non-native interface
(verbs), and the AMR code's multiple communication patterns --- multiple collectives within a
timestep, halo exchanges, and frequent load balancing --- make it highly jitter-sensitive.

\subsection{Tracing Overhead and Efficiency}
\label{sec:orca:eval-trace}

We compare ORCA with four state-of-the-art tracing tools --- TAU~\cite{Shend2006:TAU},
\scorep{}~\cite{Knupf2012:Score-P}, Caliper~\cite{Boehm2016:Caliper}, and
\dftracer{}~\cite{Devar2024:DFTracer} --- along three axes: runtime overhead, storage efficiency,
and analytics performance.

\parahead{Collection Profiles}
We set up ORCA with two tracing profiles: a lightweight profile that only captures MPI collectives,
and a detailed profile that also captures all MPI point-to-point messages and Kokkos events, except
some extremely high-volume events such as \texttt{MPI\_Test} and \texttt{MPI\_Iprobe}.
For ORCA, the lightweight profile measures the common-case overhead of always-on observability,
while the detailed profile measures the worst-case overhead of capturing all MPI and Kokkos events.
All other tracing tools are configured with an identical detailed collection profile.
Unless otherwise specified, all experiments run for 2000 timesteps, and we report the average
runtime from three runs.
Unlike tracing tools, ORCA is deployed with additional resources --- one aggregator node per 1024
ranks, plus one controller node.

% DATA: Trace size ratios for 2000 timesteps (or_trace_mpisync vs or_tracetgt)
% from tracesizes/20251229.csv
% | Ranks | mpisync (MB) | tracetgt (GB) | Ratio |
% |-------|--------------|---------------|-------|
% | 512   | ~73          | ~19.6         | 269×  |
% | 1024  | ~139         | ~34.4         | 248×  |
% | 2048  | ~266         | ~55.7         | 209×  |
% | 4096  | ~524         | ~111.6        | 213×  |
% Detailed traces are 200-270× larger than lightweight traces.

\subsubsection{Runtime Overhead}\label{sec:orca:eval-trace-runtime}
\figEvalOverhead{}
\Cref{fig:orca:eval-runtime} compares the runtime overhead of ORCA and other tracing tools for
512--4096 ranks.

\begin{itemize}[leftmargin=*]
  \item ORCA remains within 0.8--2\% overhead across all scales and profiles, while that of TAU and
        \dftracer{} degrades significantly --- from 18--20\% at 512 ranks to 56--81\% at 4096 ranks
        (\scorep{} and Caliper discussed separately below).
  \item At 4096 ranks, ORCA moves \textasciitilde550GB on-wire (\textasciitilde111GB after
        compression) over a shared fabric without QoS within 308\,s at under 2\% overhead ---
        demonstrating that RDMA-based telemetry pipelines can coexist with tightly-coupled BSP
        communication.
  \item The marginal increase in ORCA overhead between the lightweight and detailed profiles is only
        \textasciitilde0.5\%, despite collecting over 200$\times$ more telemetry (524\,MB vs
        111\,GB, respectively) --- underscoring the efficiency of Arrow serialization and
        asynchronous RDMA-based collection.
\end{itemize}

\parahead{Why ORCA Works}
TAU and \dftracer{} degrade steadily because they flush trace buffers to storage as they fill,
which produces uncoordinated flushing patterns.
This, coupled with a kernel-based data path for storage, introduces jitter that compounds at scale.
ORCA, in comparison, uses an inherently efficient format, serializes synchronously into
RDMA-registered buffers, and transports via efficient RDMA GETs.

\parahead{Caveats on \scorep{} and Caliper}
Because of difficulties running them for the full 2000 timesteps, all numbers for \scorep{} and
Caliper are extrapolated from 200 timesteps.
Both tools are designed for short characterization workloads where data is buffered in memory and
flushed once at the end.
\scorep{} additionally requires two passes --- the first to estimate buffer sizes --- and
frequently crashed with out-of-memory errors even with 64MB trace buffers.
Caliper supported periodic flushing, but the path was buggy and performance was significantly
degraded even after fixes.

\subsubsection{Storage Efficiency}\label{sec:orca:eval-trace-size}

\figEvalSizeOnly{}
\Cref{fig:orca:eval-tracesize} compares the total size, and the per-record size of the traces emitted
by ORCA and other tracers.
ORCA writes Parquet, TAU and \scorep{} are configured for OTF2, \dftracer{} writes \texttt{jsonl},
and Caliper writes its custom plaintext \texttt{.cali} format.
Key findings:
\begin{itemize}[leftmargin=*]
  \item ORCA consistently emits the smallest traces --- 3.6--3.8$\times$ smaller than TAU, and up
        to 44.3$\times$ smaller than Caliper.
  \item The differences can be partly explained by format efficiency --- OTF2, being a binary
        format, is only 1.4$\times$ larger than ORCA, while plaintext formats (\texttt{json} and
        \texttt{cali}) are 7--13$\times$ larger.
\end{itemize}

\parahead{Record Amplification}
Accounting for format efficiency, other tracing tools emit 2.5--3.3$\times$ more records than ORCA,
despite the same collection profile.
\begin{enumerate}[leftmargin=*]
  \item All tracers except \dftracer{} (and ORCA) generate separate entry and exit records for each
        traced function.
        ORCA and \dftracer{} generate one record, which contains both start time and duration.
        This is both more efficient and allows functions to be filtered by duration without requiring any
        preprocessing to pair the entry and exit records.
  \item Additionally, TAU incurs additional amplification from hardcoded derived events --- for every
        MPI send and receive, it computes a function event, a message event, and an event for
        messages received during certain phases such as \mpiw{}.
  \item Flush-induced jitter, especially with \scorep{} and Caliper, creates stragglers that other
        ranks wait on by polling --- generating additional events that inflate trace size.
\end{enumerate}

\parahead{Why ORCA Works}
ORCA's storage efficiency stems from two design choices.
First, Parquet yields an inherently compact representation through binary columnar layout,
dictionary encoding, and Snappy compression.
Second, separating derived events from collection --- computing them via OrcaFlow rather than
baking them into the tracer --- prevents record amplification from creeping in as features
accumulate.
% We discuss compression further in a later section. % TODO: add section reference

\subsubsection{Query Performance}\label{sec:orca:eval-trace-query}
\figEvalQueryOnly{}
We evaluate analytics performance using a suite of post-mortem queries (\Cref{fig:orca:eval-query}).
All queries run on traces from a smaller 20-timestep run, executed cold with caches cleared,
latencies averaged over 3 runs.

\parahead{Queries}
We use three queries of increasing complexity.
\emph{Outlier Waits} is a simple filter identifying \mpiw{} events exceeding 1\,ms.
\emph{Outlier Collectives} aggregates collective events by \texttt{swid} and filters for outliers (max
duration $>$ 10\,ms).
\emph{Timestamp Range} returns all events within the first second of execution---a query that
might power a visualization frontend, and notably not a dimension ORCA partitions on.

\parahead{Analytics Tools}
All queries execute on a single 16-core node.
\dftracer{} and Caliper use their bundled Python-based analytics libraries.
ORCA uses unmodified Polars~\cite{:Polars} with no partition awareness for simplicity--- a
handicapped configuration as the metadata from all Parquet files must be scanned for all queries.
TAU and \scorep{} are excluded as OTF2 is not a queryable format.

\parahead{Findings}
\begin{itemize}[leftmargin=*]
  \item ORCA queries complete in hundreds of milliseconds, 3--4 orders of magnitude faster than
        \dftracer{} and Caliper.
        The gap widens with scale.
  \item Query times for \dftracer{} and Caliper do not vary with query complexity.
        This is because performance is bound by parsing, not computation --- plaintext formats must be
        parsed and loaded as dataframes before any queries can be computed.
\end{itemize}

\parahead{Why ORCA Works}
While all analytics-oriented tracers ultimately load data into dataframes, they persist traces as
plaintext formats requiring expensive ETL to materialize dataframes before any query executes.
ORCA preserves dataframe structure end-to-end, persisting directly as Parquet --- enabling zero-ETL
post-hoc analytics.
Additionally, the layout confers two kinds of pruning: file-level pruning via (timestep, rank)
partitioning, and intra-file pruning via rowgroup statistics.
Our benchmark does not exploit the former --- unmodified Polars with lazy execution must scan all
file footers.
Adopting Parquet-based lakehouse metadata formats~\cite{:Iceberg,Armbr2020:Deltalake} would add a
catalog layer over partitioned Parquet, making partition pruning automatically available via
standard query engines~\cite{:Polars,Raasv2019:DuckDB}.
Even without partition pruning, ORCA outperforms by 3--4 orders of magnitude.
Conventional formats offer no such pruning --- as traces grow, larger analytics clusters must be
provisioned for ETL before queries can begin.

\subsubsection{Tradeoff Analysis}\label{sec:orca:eval-trace-tradeoffs}

\figEvalArch{}
We isolate the impact of key architectural decisions using a 512-rank run (\Cref{fig:orca:eval-arch}).

\parahead{Compression}
\dftracer{} supports gzip compression, which reduces its trace size from $12.2\times$ that of ORCA
to $1.6\times$ (\cref{fig:orca:eval-arch-dft}), but at the cost of increasing runtime overhead from
18\% to 504\%.
Compression on MPI ranks is unviable because it competes with the application for CPU, amplifying
jitter.
ORCA offloads compression to dedicated aggregator cores, keeping it off the critical path.
Additionally, compressed jsonl cannot be queried without full decompression, while ORCA's causal
partitioning and Parquet's rowgroup structure reduce decompression overhead.

\parahead{Transport Layer}
Switching ORCA from \texttt{verbs} to the libfabric \texttt{TCP} provider increases overhead from
1--2\% to 45--55\% (\cref{fig:orca:eval-arch-tcp}).
Even with the same Arrow-based data path, TCP suffers because retrieval incurs CPU overhead on
simulation ranks, introducing jitter.
This gap reflects the complementary nature of ORCA's transport choices: Arrow enables efficient
serialization into contiguous RDMA-friendly buffers, and RDMA allows asynchronous retrieval with
minimal CPU overhead while the application runs undisturbed.
These results underscore the cost of TCP-based transport for tightly-coupled BSP workloads.
GPU-based codes may tolerate the higher transport overhead, but event-based streaming tools still
lack the columnar, causally-partitioned data path that enables ORCA's analytical efficiency.


\subsection{ORCA-Native In-Situ Workflows}
\label{sec:orca:eval-insitu}

\tabEvalFlows{}

\figEvalOrcaNative{}

While ORCA can enable conventional tracing workflows, it is fundamentally designed for in-situ
telemetry analytics --- directly materializing views from a running application.
Plan operators can be pushed all the way to the MPI tier, leveraging the application's own
resources for data reduction at source.
We now evaluate these capabilities, measuring the runtime overhead of operator pushdown, and its
effectiveness at reducing captured data volume.
% The previous section compared ORCA against existing tools on their own terms. We now evaluate
% capabilities that have no analog in existing tooling: real-time in-situ telemetry analytics with
% operator pushdown.

\parahead{Workflows}
\Cref{tab:orca:ontv-flows} shows three flows of increasing complexity, designed for AMR codes.
The first flow demonstrates a monitoring workflow: measuring the volume of \texttt{MPI\_Wait}
exceeding 1ms, while the second flow persists all outlier events to Parquet for offline analysis.
The third flow tracks compute load imbalance in AMR codes --- it filters for known compute kernels,
computes duration aggregates by rank, and streams details to Parquet and percentiles to DuckDB for
real-time dashboards.


\parahead{Observations}
\Cref{fig:orca:eval-ontv} shows the runtime overhead and the data reduction at \tiermpi{} for the three flows.
\begin{itemize}[leftmargin=*]
  \item
        All workflows incur an overhead of $\le 2\%$ across all scales despite executing pushed-down
        operators on the MPI tier of a CPU simulation, illustrating the efficiency of Arrow-based analytics.
        Computation budget for pushdown only grows with modern GPU codes.
  \item Pushdown reduces transmitted data volume by $98.5\%-99.999\%$.
        Leveraging MPI nodes for analytics reduces fabric utilization, residual work at upper tiers, and
        shortens the feedback loop.
        The \emph{capture-everything} nature of traditional tracing is redundant for monitoring and
        debugging---a small fraction of aggregated and filtered data suffices to materialize the required
        views.
        %   only a small amount of aggregated data is required to construct the views required for
        %   these workflows.
  \item The \mpiw{} flows do not sample---they perform continuous predicate evaluation for every
        single \mpiw{}, emitting only outliers.
        Comprehensive continuous observability of this type is impractical without pushdown.
\end{itemize}

\subsection{Closing the Loop: Agentic Debugging}\label{sec:orca:eval-agent}

While ORCA's control plane was designed for human operators, the emergence of LLM-based coding
agents---capable of issuing commands and analyzing outputs---offers a way to validate the
end-to-end feedback loop.
To evaluate whether ORCA's interfaces are expressive enough for autonomous debugging, we use an LLM
agent as a proxy for a human operator---tasking it with localizing the anomaly from
\Cref{fig:orca:bg-volatile}, the same bug that took months to localize manually.

\parahead{Setup}
The anomaly is reproduced within the same AMR code used in previous experiments at 4096 ranks,
using a placement policy that exacerbates the anomaly's frequency.
The code is instrumented with a hierarchy of Kokkos region annotations, with additional annotations
in functions of interest added by us.
Claude Code (Opus 4.5)~\cite{:ClaudeCode} ran on the head node with access to the ORCA
controller via CLI and the codebase.
The agent was tasked with monitoring collectives for outliers and progressively collecting only the
Kokkos and MPI events necessary to localize the problem.
The LLM issued all ORCA commands and performed all data analysis using off-the-shelf
Polars~\cite{:Polars} on the head node.

% The anomaly is reproduced within the same AMR code used in previous experiments at 4096 ranks,
% albeit with a different placement policy that exacerbates the anomaly's frequency.
% The code is instrumented with a hierarchy of Kokkos region annotations, with additional annotations
% in functions of interest added by us.
% An LLM agent ran on the head node with access to the ORCA controller via CLI and the codebase.

% The agent was tasked with monitoring collectives for outliers and progressively collecting only the
% Kokkos and MPI events necessary to localize the problem.
% Of the three LLM agents tested---Claude Code~\cite{:ClaudeCode} with Opus 4.5, OpenAI
% Codex~\cite{:CodexCLI} with GPT 5.2 Thinking, and Gemini CLI~\cite{:GeminiCLI} with Gemini 3---only
% Claude Code completed the exercise with some human guidance.
% The LLM issued all ORCA commands and performed all data analysis using off-the-shelf
% Polars~\cite{:Polars} on the head node.

\figEvalAgentic{}

\parahead{Timeline}
\Cref{fig:orca:eval-agentic} shows the debugging session: user and LLM interactions at top,
per-timestep event volume at bottom.
After a warmup phase, the LLM began the localization exercise at 272\,s by monitoring collectives
for outliers (\flowmark{1}).

\begin{enumerate}[leftmargin=*]
  \item At 384\,s, the LLM analyzed collective percentiles and identified outlier stragglers.
        It installed a flow to collect Kokkos events with duration $> 1$\,ms (\flowmark{2}).
  \item At 510\,s, it began investigating collected Kokkos events, correlated outliers with
        the collective pattern, and attributed
        the anomaly to a \texttt{map.clear()} call during AMR redistribution.
  \item At 1042\,s, the user rejected the \texttt{map.clear()} hypothesis as an inadequate explanation for latencies between 50--100\,ms.
        The LLM investigated further and identified that \texttt{map.clear()} triggered a buffer destructor
        calling \mpiw{}.
  \item At 1191\,s, the user asked it to prove the \mpiw{} hypothesis using ORCA.
        The LLM proposed capturing all \mpiw{} events, but the user rejected this and asked for a
        duration-based filter.
        % for volume.
        A revised flow with a 1\,ms filter was accepted (\flowmark{3}).
        %   It revised the flow with a 1\,ms filter, which was accepted (\flowmark{3}).
  \item At 1335\,s, the LLM confirmed outlier waits as the root cause.
        Asked to trace the control flow leading to the anomalous \mpiw{}, it generated a report at 1605\,s.
        The application then resumed with baseline collective monitoring (\flowmark{1}).
\end{enumerate}

\parahead{Observations}
% The exercise completed in 27 minutes, with the LLM driving most decisions.
% Reverting to baseline monitoring afterwards reduced data volume by 144$\times$.
% While this exercise stopped at \mpiw{}, support for dynamic instrumentation (eBPF, DynInst) would
% enable the investigation to continue into the MPI library and fabric layers.
% We had earlier traced this spike to misplaced \texttt{ACKs} within the fabric driver, triggering an
% expensive timeout/retransmit path.
% It was subsequently fixed with an asynchronous drain queue for stuck requests.
% Beyond localization, ORCA supports always-on production monitoring of mitigations: the same
% signature that identified the anomaly can be deployed as a metric, enabling regressions to be
% detected immediately.
The exercise validates ORCA's closed-loop capabilities.
The agent autonomously correlated high-level telemetry with code, requested targeted detail, and
traced the root cause to where instrumentation ended --- completing in 27 minutes what took months
of manual investigation.
Reverting to baseline monitoring afterwards reduced data volume by 144$\times$.
User interventions were minor: rejecting an incomplete but accurate explanation and asking for
basic duration-based filtering.
% Some additional annotations in the targeted regions made the exercise easier
The exercise was made easier with some additional instrumentation in relevant regions, but the
result still demonstrates that ORCA's feedback loop enables online localization workflows as long
as the right instrumentation is present.
Support for dynamic instrumentation~\cite{:eBPF, Berna2011:DynInst} would close this gap, enabling
investigations to continue into the MPI and fabric layers.

Manual investigation had previously traced the anomalous spikes to misplaced \texttt{ACKs} within
the fabric driver, triggering an expensive timeout/retransmit path --- subsequently fixed with an
asynchronous drain queue.
Beyond localization, ORCA supports always-on production monitoring of mitigations: the same
signature that identified the anomaly can be deployed as a metric, enabling regressions to be
detected immediately.
